\section{Conclusion} \label{sec:conclusion}
% \todo{ER: Recapt contributions}
\vspace{-5pt}
To enable realistic, reproducible, and robust evaluation for machine learning on temporal graphs, we present the Temporal Graph Benchmark, a collection of challenging and diverse datasets. \name datasets are diverse in their dataset properties as well as being orders of magnitude larger than existing ones. \name includes both \emph{dynamic link property prediction} and \emph{dynamic node property prediction} tasks, while providing an automated pipeline for researchers to evaluate novel methods and compare them on the \name leaderboards. 
In dynamic link property prediction, we find that model rankings can vary significantly across datasets, thus demonstrating the necessity to evaluate on the diverse range of \name datasets. Surprisingly for dynamic node property prediction, simple heuristics such as persistence forecast and moving average outperforms SOTA methods such as TGN. This motivates the development of more TG methods for node-level tasks.

\textbf{Impact on Temporal Graph Learning.}
Significant advancements in machine learning are often accelerated by the availability of public and well-curated datasets such as ImageNet~\cite{deng2009imagenet} and OGB~\cite{hu2020open}. We expect \name to be a common and standard benchmark for temporal graph learning, helping to facilitate novel methodological changes. 



\textbf{Potential Negative Impact.}
If \name becomes a widely-used benchmark for temporal graph learning, it is possible that future papers might focus on \name datasets and tasks, which may limit the use of other TG tasks and datasets for benchmarking. To avoid this issue, we plan to update \name regularly with community feedback as well as adding additional datasets and tasks. 


\textbf{Limitations.}
Firstly, \name only considers the most common TG evaluation setting as discussed in Section~\ref{sec:task}, namely the \emph{streaming} setting. We also discuss other possible settings in Appendix~\ref{app:eval}.
Depending on a specific application, a different setting might be more suitable (such as forbidding test time node updates). Secondly, \name currently only contains datasets from five domains, while many other domains such as biological networks are not included. We plan to continue adding datasets to \name to further increase the dataset diversity in \name. 
